\documentclass{article}
\usepackage{amsmath,amssymb,amsthm}
\newtheorem{theorem}{Theorem}[section]
\newtheorem{lemma}[theorem]{Lemma}
\theoremstyle{definition}
\newtheorem{definition}[theorem]{Definition}

\title{Sums of Squares}
\author{A. Student}
\date{October 2026}

\begin{document}
\maketitle
\begin{abstract}
A short article: sections, numbered equations, theorems with proofs, and cross-references.
\end{abstract}
\tableofcontents

\section{Introduction}\label{sec:intro}
Which integers are sums of two squares? We collect the classical answer; see Theorem~\ref{thm:fermat}
and Section~\ref{sec:proof}.

\section{Statement}
\begin{definition}
An integer $n \ge 0$ is a \emph{sum of two squares} if $n = a^2 + b^2$ for some $a, b \in \mathbb{Z}$.
\end{definition}

\begin{lemma}\label{lem:product}
Products of sums of two squares are sums of two squares:
\begin{equation}\label{eq:brahmagupta}
(a^2+b^2)(c^2+d^2) = (ac-bd)^2 + (ad+bc)^2 .
\end{equation}
\end{lemma}

\begin{theorem}[Fermat]\label{thm:fermat}
An odd prime $p$ is a sum of two squares if and only if $p \equiv 1 \pmod 4$.
\end{theorem}

\section{Proof sketch}\label{sec:proof}
\begin{proof}
If $p = a^2 + b^2$, then $p \equiv 1 \pmod 4$, since squares are $0$ or $1$ modulo $4$. Conversely,
choose $x$ with $x^2 \equiv -1 \pmod p$ and apply Thue's lemma; by \eqref{eq:brahmagupta} the rest follows.
\end{proof}

\begin{align}
  \sum_{n=1}^{\infty} \frac{1}{n^2} &= \frac{\pi^2}{6}, &
  \int_0^1 x^n \, dx &= \frac{1}{n+1}.
\end{align}
\end{document}
